\documentclass[11pt,a4paper]{ctexart}
\usepackage[margin=2.4cm]{geometry}
\usepackage{amsmath,amssymb}
\usepackage{fancyhdr}
\usepackage{enumitem}
\linespread{1.18}
\setlength{\headheight}{15pt}
\setlength{\parskip}{0pt}

\pagestyle{fancy}
\fancyhf{}
\fancyhead[L]{\small Team 2209466}
\fancyhead[R]{\small \thepage}
\renewcommand{\headrulewidth}{0.4pt}
\fancypagestyle{plain}{\fancyhf{}\renewcommand{\headrulewidth}{0pt}}

\ctexset{
  section/format = \Large\bfseries\raggedright,
  subsection/format = \large\bfseries\raggedright,
  subsubsection/format = \normalsize\bfseries\raggedright
}
\setcounter{secnumdepth}{3}

\begin{document}

%====================== 摘要页 ======================
\thispagestyle{empty}
\begin{center}
{\small 2022 MCM/ICM\quad 选题 A\quad 队伍控制号 2209466\quad 摘要页（Summary Sheet）}\\[18pt]
{\LARGE\bfseries 自行车运动员的最优功率分配}\\[8pt]
{\large （Optimal Power Distribution of a Cyclist）}\\[10pt]
\end{center}

\begin{center}
{\large\bfseries 摘\quad 要}
\end{center}

\vspace{6pt}
自行车公路赛一直以来都备受关注且竞争激烈。对于自行车运动员而言，在赛道上合理地分配功率至关重要。为了帮助骑行者合理地分配其有限的体力，我们建立了一个模拟骑行者功率输出的模型，以寻找其功率的最优分配方案。

首先，我们建立了一个疲劳生理模型（Fatigue Physiological Model）来刻画骑行者的功率曲线，该曲线描述了骑行者的平均输出功率随持续输出时间变化的函数关系。我们从有氧呼吸与无氧呼吸的角度出发，充分考虑能量的有限性以及能量恢复的过程，建立了有氧与无氧能量疲劳模型（Aerobic and Anaerobic Energy Fatigue Model）。我们使用普通最小二乘法（ordinary least squares）将模型与所获得的数据进行拟合，以验证模型的有效性，并考虑了不同运动员类型与性别的情形。我们给出了不同骑行者的功率曲线，并探究了其各项特征与模型参数之间的关系。具体结果见图 6 及其分析。

其次，我们建立了一个个体最优功率分配模型（Optimal Individual Power Distribution Model），以给出最优的功率分配策略，该模型描述了骑行者所施加的功率与赛道上位置之间的关系。我们建立了一个单目标规划模型，以比赛用时为优化目标，以骑行者状态与赛道条件为约束条件，以骑行者所施加的功率为决策变量。我们构建了一个阻力模型来描述阻力，并使用赛道起伏模型与弯道模型来描述赛道条件。在上述模型与骑行者功率曲线的基础上，我们采用动态规划（dynamic programming）求得最优功率分配策略及相应的最短用时。我们还考察了赛道条件与骑行者状态对骑行者功率分配的影响。功率输出的分配图与相关讨论将在第 5.4 节中给出。

第三，我们针对天气条件进行了灵敏度分析，其中包括风向与风速。我们将先前的阻力模型改进为风阻模型。在由 8 个风向与 3 个风速等级构成的 24 种环境中，我们计算了最短比赛用时，结果见图 13。完赛时间变化很大，最大达到 16\%，因此我们认为结果对天气条件是敏感的。

第四，我们对骑行者偏离功率分配策略的情形进行了灵敏度分析。我们使用 MATLAB 对最优功率分配策略进行随机扰动，并以骑行者实际输出功率与各位置上最优功率之间偏差绝对值的平均值，作为偏离程度的度量。完赛时间与偏离程度之间的关系见表 4。直到平均偏差变化达到 5\% 时，完赛时间才增加 10\%。因此，我们认为模型结果对骑行者偏离目标功率的情形是不敏感的。对于接近最优的完赛时间，赛道上各个位置的时间参考区间见表 5。

最后，我们考虑了团队计时赛的特征，将个体比赛模型推广为团队策略决策模型（Team Strategy Decision Model）。我们考虑了队员之间的挡风效应，修正了风阻模型，并提出了轮换领骑策略（draft strategy）与掉队策略（drop strategy）。此外，我们还为个人计时赛提供了一份指导。

\vspace{6pt}
\noindent\textbf{关键词：}个人计时赛；功率曲线；最优功率分配；动态规划

\newpage
\pagestyle{fancy}

%%====================== 正文 ======================
\section{引言}

\subsection{问题背景}

公路自行车比赛种类繁多，其中个人计时赛要求自行车运动员独自骑行一段距离，以完赛时间作为成绩，用时最短者获胜。

输出功率与其持续时间之间的曲线关系可以作为高强度运动的一项本质特征，这一概念被称为功率曲线（power curve）。从本质上讲，该概念描述的是高强度运动可耐受的持续时间。它表征了运动员在某一特定时间段内所能输出的最大功率，在优化运动训练方案与运动表现、以及改善慢性疾病患者的生活质量等方面具有巨大的应用潜力[1]。

功率曲线的形态本质上由人体三种供能系统所占比例随时间的迁移所决定。在数秒至十余秒的极短时段内，磷酸原系统与无氧糖酵解系统占据主导，此时运动员所能输出的功率可达其有氧功率的数倍；随着持续时间延长，无氧储备逐渐被消耗，功率随之迅速下降；当持续时间延长到数分钟乃至数十分钟之后，功率的下降趋于平缓，并逼近一个几乎不再随时间下降的平台，该平台的水平由有氧代谢能力所决定。因此，若以持续时间的对数作横坐标、以功率作纵坐标，功率曲线通常呈现先陡降、后趋缓的形态，其转折点的位置与形状同时取决于运动员的无氧储量与有氧功率水平。也正因为如此，功率曲线可以被视为运动员能量储备的一张“身份证”：曲线的短时端刻画爆发能力，长时端刻画耐力水平，二者的相对高低则决定了运动员的类型归属。

不同类型的骑行者的功率曲线形状各不相同。在给定的一次计时赛骑行中，他们需要管理自己的功率分配，以便以最少的时间完成比赛。图 1 展示了一位正在比赛中的自行车运动员。

\subsection{问题的重述}

考虑到问题陈述中所给出的背景信息与限制条件，我们需要建立一个模型来确定一名骑行者在一条赛道上的最优功率分配，并解决以下问题：

\begin{itemize}
\item \textbf{问题 1}：给出两类骑行者的功率曲线（power profile）的定义，并考虑不同性别的骑行者。其中一名骑行者是计时赛专家（time trial specialist）；另一名则要求属于另一种类型。
\item \textbf{问题 2}：将上述模型应用于三条具体的计时赛赛道。
\item \textbf{问题 3}：分析天气条件对模型结果的影响，以及模型对该影响的灵敏度。
\item \textbf{问题 4}：分析模型对偏离所得功率分配结果的灵敏度。
\item \textbf{问题 5}：讨论如何发展该模型，以便为参加团队计时赛的六人团队提供最佳的功率使用方案。
\end{itemize}

上述五个问题构成了一个由个体到团队、由确定性到灵敏度的递进结构：问题 1 要求刻画骑行者的能力边界，它是后续所有优化问题的输入；问题 2 要求在给定能力边界与赛道条件下最优地使用能量；问题 3 讨论外部环境扰动对最优策略的影响；问题 4 讨论骑行者无法精确执行策略时的鲁棒性；问题 5 则将单主体模型扩展为多主体协同问题。与这一结构相对应，我们将建模工作划分为四个层次：能力模型（疲劳生理模型）、环境与赛道模型（阻力模型与赛道抽象模型）、优化模型（个体最优功率分配模型）以及推广模型（团队策略决策模型）。本文的组织顺序亦遵循这一层次，先解决“骑行者能做什么”，再解决“骑行者应当做什么”，最后讨论“团队应当如何协同”。

\subsection{文献综述}

本文所针对的问题是骑行者在骑行过程中的功率分配，此问题在以往的研究中已有讨论。Andrew M. Jones 等人将功率曲线引入间歇性高强度运动，并使用已有的双参数模型（功率与持续时间之间的双曲关系）对运动员的疲劳监测作出了更好的预测[3]。Morton 将双参数模型扩展为三参数模型，并建立了一个新的功率曲线模型[4]。Andrew Coggan 博士等人使用三参数模型，通过统计数据分析了不同类型骑行者的优势与劣势，并得到了相应的功率曲线[5]。

关于如何为骑行者合理地分配功率以赢得比赛，同样有大量研究。Monod 与 Scherrer 提出的临界功率模型（critical power model）是被广泛认可的、用于优化配速策略与分析骑行者耐力能力的模型。它由两部分构成：无限的有氧部分与有限的无氧部分[6]。更为复杂的模型是三分量 Morton--Margaria 模型，或称三罐模型（3-Tank model）[7]，它是另一种可行的生物能量模型，为自行车比赛中的功率分配研究提供了方向。

\subsection{我们的工作}

我们需要一个适用于任何类型骑行者的模型，用以刻画骑行者在赛道上的位置与其所施加的力之间的关系。我们的工作主要包括以下几个部分：

首先，我们参考以往的文献，在三罐模型的基础上提出了我们的自行车运动员疲劳生理模型（Cyclist's Physiological Model of Fatigue），其中包括有氧能量疲劳模型与无氧能量疲劳模型。我们使用该模型来解释骑行者的功率曲线。此后，我们对不同骑行者的曲线进行了最小二乘拟合，并使用所得到的参数将骑行者区分为不同的类别。

接下来，我们建立了一个以空气动力阻力与滚动摩擦阻力为主导的阻力模型，并在考虑天气条件的影响时对该阻力模型进行了修正。此外，我们将赛道抽象为作为比赛进程函数的地面起伏，以及作为比赛进程函数的弯道。这就是我们的个体最优功率分配模型（Optimal Individual Power Distribution model）。此后，我们使用具体的赛道数据与所提出的疲劳生理模型，通过动态规划对其求解，得到了骑行者在赛道上的最佳功率—比赛进程曲线以及最佳完赛时间。

然后，我们分别针对风速与风向、以及功率偏差进行了灵敏度分析，并分析了模型的优点与缺点。

最后，我们将模型推广到团队计时赛。为了考虑队伍之间的阻力减小效应，我们重新构建了团队比赛中个体的阻力模型。此后，我们考虑了团队比赛中可能采用的策略——轮换领骑（Draft）与掉队（Drop），并将这些策略发生的位置作为新的决策变量，从而在理论上可以获得团队比赛中的最优解。

为了便于对整体模型的理解，我们给出了一张流程图来描述我们的工作流程。

\section{假设与理由}

考虑到实际问题总是包含许多复杂的因素，首先，我们需要作出合理的假设以简化模型，而每一条假设之后都紧随其相应的解释：

\begin{itemize}
\item \textbf{假设 1}：骑行者在比赛过程中绝不会受到他人的影响。根据国际自行车联盟（UCI）的自行车个人计时赛规则，骑行者按一定的时间间隔单独出发。因此，骑行者在比赛中不会受到他人的影响。
\item \textbf{假设 2}：骑行者在比赛过程中的能量是有限的。这一方面来自题目的限定，另一方面，骑行者在实际比赛中通常不会补充能量。
\item \textbf{假设 3}：阻力的唯一来源是空气的黏性阻力与地面的滚动摩擦。因为这两种阻力占了最大的比重[8]，而来自自行车链条传动效率的微小阻力可以忽略不计。
\item \textbf{假设 4}：骑行者沿赛道作一维运动，不考虑赛道宽度。这是对模型的一个合理简化。对横向运动的分析会引入不必要的复杂性。
\item \textbf{假设 5}：天气对运动员的影响只存在于风向与风速之中，不考虑降雨或降雪的影响。这是对模型的一个合理简化，因为实际比赛一般不会在更为恶劣的天气条件下进行。
\end{itemize}

上述五条假设分别对应模型中可以被忽略或被合并的因素，其共同的目的是在不损害主要矛盾的前提下压缩模型的自由度。假设 1 与假设 4 保证了骑行者的运动可以作为一个单质点的一维运动来处理，从而使得赛道上任意位置的状态只由该位置本身决定，而不依赖于横向位置或其他选手；假设 2 与假设 3 界定了能量的输入与输出边界，使得功率曲线的拟合与后续的可行性判定有了明确的口径；假设 5 则界定了外部扰动的维数，使灵敏度分析可以在一个二维参数空间（风向与风速）中完成。需要强调的是，上述假设并非对现实的简单回避，而是在后续的灵敏度分析中逐一放开或被修正：例如在天气条件分析中，我们通过修正阻力模型而部分放开了假设 5；在团队计时赛的推广中，我们通过引入位置相关的阻力系数而部分放开了假设 1。

\section{定义与符号}

\subsection{定义}

\textbf{功率体重比（Power-to-weight ratio）}：骑行者输出功率与其自身体重之比。该值越高，爬坡能力越强。

\textbf{功能阈功率（Functional Threshold Power，FTP）}：骑行者在全力且稳定骑行一小时内所能达到的最大平均功率。它是骑行者有氧能力的一项重要指标。它又被称为乳酸阈，通常通过测试确定。

\textbf{临界功率（Critical Power）}：骑行者在给定时间段内以全力且稳定的方式骑行所能获得的最大平均功率。

\textbf{急弯（Sharp turn）}：在 50 米范围内转角大于或等于 90 度的弯道。

上述四个定义构成了本文与题目之间、以及不同模型层次之间的共同语言。功率体重比刻画的是单位质量所能获得的推进能力，因而与坡度的影响直接相关；功能阈功率刻画的是有氧供能的上界，它在我们的模型中既是无氧呼吸的起始功率，也是长时间骑行时功率曲线的渐近水平；临界功率刻画的是某一给定时间段内的最大稳定输出，用于与功能阈功率相互对照；急弯的定义则把一个连续的几何概念离散化为一个可以在赛道上判定的阈值条件，从而使得弯道可以被处理为速度检查点。凡此种种，均是为了把生理学与几何学的概念转化为模型中可以计算的量。

\subsection{符号}

本文所使用的主要数学符号列举如下（依次为符号、说明与单位）：

\begin{itemize}
\item $P_A$：有氧呼吸的能量输出功率，单位 W/kg；
\item $r_A$：有氧池能量比例的消耗速率，单位 s$^{-1}$；
\item $P_{An}$：无氧呼吸的能量输出功率，单位 W/kg；
\item $r_{An}$：无氧池能量比例的消耗速率，单位 s$^{-1}$；
\item $W'$：无氧能量储量，单位 J；
\item $t$：持续时长，单位 s；
\item $k$：无氧池能量比例的恢复速率，单位 s$^{-1}$；
\item $P$：骑行者的总输出功率，单位 W/kg；
\item $F_{\mathrm{res}}$：总阻力，单位 N；
\item $v_{ra}$：骑行者相对于空气的速度，单位 m/s；
\item $\mu$：滚动摩擦系数，无量纲；
\item $F$：骑行者所提供的动力，单位 N；
\item $x$：骑行者在赛道上的位置，单位 m；
\item $\theta$：赛道的坡度角，单位度；
\item $h$：海拔高度，单位 m；
\item $\mathcal{C}$：赛道的弯道集合，无量纲；
\item $\varphi$：行进方向与风向之间的夹角，单位度；
\item $v_w$：风速，单位 m/s；
\item $D(x)$：轮换领骑的最优决策函数，无量纲。
\end{itemize}

注：还有一些变量未在此列出，将在各节中详细讨论。

在进入建模之前，需要说明本文的归一化约定：有氧池与无氧池的能量储量均被归一化为一，因而二者的剩余量始终落在区间 $[0,1]$ 上；功率一律按骑行者体重归一化，单位为 W/kg，这既消除了体重对功率绝对值的影响，也使得不同性别、不同类型骑行者的曲线可以直接比较；位置 $x$ 以米为单位沿赛道弧长计量，因而赛道的各项几何属性均可表示为 $x$ 的函数。凡涉及时间的量，若不特别说明，均以秒为单位。

\section{自行车运动员的疲劳生理模型}

下面我们详细阐述我们的模型。在三罐模型的基础上，我们提出了一个有氧与无氧能量疲劳模型，并使用最小二乘拟合曲线来获得不同类型运动员的功率—时间曲线。

\subsection{模型的建立}

\subsubsection{功率曲线模型}

上文提到的三罐模型[7]在理论上认为有氧来源可以提供无限的能量。然而在实际过程中，根据假设 2，总能量是有限的。因此，我们可以假定有氧呼吸功率与持续时长之间存在负相关关系。

当一个人处于全力活动的状态时，由于有氧呼吸的速率有限，有氧呼吸所提供的功率是有限的。当输出功率超过功能阈功率（FTP）时，就需要无氧呼吸来补足。因此，有氧呼吸功率属于一个具有上界的区间，其消耗速率也存在上界。

由上述两点，结合其他能量模型，我们将有氧池的能量归一化，并定义骑行者消耗有氧池能量的速率为：

\begin{equation}
r_A = \frac{\mathrm{d}S_A}{\mathrm{d}t} = \frac{P_A - a}{b},
\end{equation}

其中，$P_A$ 为有氧呼吸的能量输出功率；$a$、$b$ 均为待定参数。

该式的物理含义是：有氧池的能量消耗速率与当前的有氧输出功率 $P_A$ 成正比，其比例由参数 $b$ 控制，而 $a$ 则刻画出有氧供能的下限功率水平，即当 $P_A$ 降至 $a$ 以下时，有氧池不但不再被消耗，反而开始积累。由于有氧池的能量已被归一化，$r_A$ 的量纲为 s$^{-1}$，它与输出功率的关系构成了“功率越高、可维持时间越短”这一现象在能量层面的解释。参数 $a$ 与 $b$ 由后文的拟合确定，它们因运动员的有氧能力而异：有氧能力越强的运动员，其 $a$ 越大、$b$ 越小，因而在相同的输出功率下消耗得更慢。

当输出功率超过功能阈功率时，有氧呼吸功率达到饱和，超出的部分需要由无氧呼吸来补充。并且无氧呼吸的初始功率可以近似为 FTP。

由双参数模型可知，无氧呼吸功率与持续时长之间存在双曲关系，即无氧能量池的能量储量是一个定值。因此，无氧呼吸部分的输出功率可以表示为：

\begin{equation}
P_{An} = \frac{W'}{t},
\end{equation}

其中，$W'$ 表示无氧能量池的能量储量；$t$ 为持续时长。

该式即双参数模型在这一问题中的具体形态：若运动员仅依靠无氧能量池以恒定功率输出，则其可维持的时间与功率的乘积恒等于储量 $W'$，因而功率与时间在双曲线上互为倒数关系。这一关系解释了功率曲线在短时端的陡峭形态——时间缩短一半，所需的功率就要翻倍。

考虑到人体的输出功率越高，能量的消耗就越快。根据双参数模型，我们将无氧能量池归一化为一，并定义骑行者消耗无氧能量池能量的速率为：

\begin{equation}
r_{An} = \frac{\mathrm{d}S_{An}}{\mathrm{d}t} = \frac{P_{An}}{\lambda P_A},
\end{equation}

其中，$\lambda$ 为人体无氧能量消耗速率与有氧呼吸输出功率之比。

式中 $\lambda$ 的作用是把无量纲的归一化储量与具有量纲的功率联系起来：它表示单位有氧输出功率所对应的无氧消耗速率，因而反映了运动员在超出 FTP 之后“透支”无氧储备的猛烈程度。$\lambda$ 越小，说明同样的超出功率所带来的无氧消耗越慢，运动员在高功率区间的耐受能力越强。

此外，骑行者可能选择暂时超出其功率曲线的限制。然而，在输出功率降至 FTP 以下之后，骑行者需要额外的时间以较低的功率水平来恢复无氧能量池的能量，我们据此定义骑行者恢复无氧能量池能量比例的速率：

\begin{equation}
\frac{\mathrm{d}S_{An}}{\mathrm{d}t} = -\frac{k\left(\mathrm{FTP} - P\right)}{W'}, \qquad P < \mathrm{FTP},
\end{equation}

其中，$k$ 为恢复系数，它意味着当输出功率超出曲线的限制时，骑行者大约需要两倍的时间来恢复能量池中所储存的能量。

该式中的负号表示恢复过程使无氧池的能量比例上升，其速率与当前的功率缺口 $\mathrm{FTP} - P$ 成正比，并与储量 $W'$ 成反比，这一结构保证了量纲的一致性（s$^{-1}$）。取 $k \approx 0.5$ 的依据在于：恢复过程本身同样需要消耗有氧能量，且乳酸等代谢产物的清除并非瞬时完成，因而骑行者在低功率区间恢复无氧储备所需的时间，约为其在高功率区间消耗同等储备所需时间的两倍。这一“消耗快、恢复慢”的不对称性，正是最优功率分配策略中出现“冲刺—恢复”交替波动的根本原因。

因此，当输出功率保持恒定时，骑行者维持该功率的时间可以表示为：

\begin{equation}
t = \frac{W'}{P - \mathrm{FTP}}, \qquad P > \mathrm{FTP}.
\end{equation}

由该式可以看出，当 $P$ 仅略高于 FTP 时，可维持时间趋于无穷，这符合有氧供能可以长期持续的生理事实；而当 $P$ 远高于 FTP 时，可维持时间迅速缩短，与之对应的是短时极限功率的出现。

需要说明的是，上述公式并未涉及无氧池能量的恢复。这是因为骑行者保持功率恒定。当输出功率高于 FTP 时，无氧池的能量被消耗而不会恢复；当输出功率低于 FTP 时，无氧池的能量不会被消耗。

综上所述，我们可以得到总输出功率与持续时长之间的关系为：

\begin{equation}
P(t) = \mathrm{FTP} + \frac{W'}{t}\left(1 - \mathrm{e}^{-\gamma t}\right) + \frac{\alpha}{t + \beta},
\end{equation}

其中，$\alpha$、$\beta$ 为调整系数，用以保证该函数满足 $\lim\limits_{t \to 0} P(t) = +\infty$ 且 $\lim\limits_{t \to \infty} P(t) = \mathrm{FTP}$；$\gamma$ 表示有氧呼吸供能时的能量消耗速率。

该式的三项分别对应三类贡献：常数项 FTP 是有氧供能的上界，也是长时间骑行的渐近功率；第二项是无氧供能项，它随时间的延长而以 $1/t$ 的方式衰减，其中因子 $(1-\mathrm{e}^{-\gamma t})$ 用于修正极短时间内的供能速率，使模型不至于在 $t \to 0$ 时给出无穷大的功率；第三项则是对有氧供能随能量消耗而衰减的修正，它在短时间内抬高了可输出功率，而在长时间内趋于零。三者叠加，使得曲线在短时端足够陡峭、在长时端平滑地趋于 FTP，从而同时刻画了运动员的爆发能力与耐力水平。

\subsection{结果与讨论}

\subsubsection{模型验证}

上述分析为各类自行车运动员的功率—时间曲线提供了一个基础模型。不同的自行车运动员由于其自身条件不同，参数也会有所差异，从而影响其功率—时间曲线。对于不同的运动员，可以通过测量获得具体的功率—时间曲线以反映各项指标，进而得到相应的有氧能量池与无氧能量池的消耗函数，这对疲劳监测等方面具有重要意义。

利用上文提到的数据，我们采用最小二乘法对曲线进行拟合，功率—时间曲线的拟合结果如下：

具体而言，对每一名（或每一类）骑行者的功率—时间数据点 $\{(t_i, P_i)\}_{i=1}^{n}$，我们以第 (6) 式为模型函数，构造残差平方和
\[
Q\left(\mathrm{FTP}, W', \gamma, \alpha, \beta\right) = \sum_{i=1}^{n}\left[P_i - P\left(t_i\right)\right]^2 ,
\]
并通过求解 $\partial Q / \partial \theta = 0$（$\theta$ 为待估参数向量）得到参数的最小二乘估计。由于模型关于参数是非线性的，实际计算中采用带边界的非线性最小二乘迭代，并以 FTP 的生理学先验范围作为约束，以避免迭代收敛到无生理意义的参数值。拟合的优度以决定系数与残差的标准差共同衡量。

首先需要说明的是，由于人类在功率输出上存在极限，因此输出时间极短时的输出功率不予考虑。所以图像的横坐标从 10 s 开始。为了更公平地描述骑行者的功率曲线，我们以骑行者所施加的、按体重归一化后的功率作为纵坐标。

由上述拟合结果来看，从拟合效果上看，两次拟合得到的结果曲线与原始数据吻合良好，尤其是在 $t$ 较大的区段。女子的结果曲线比男子的结果曲线更贴合原始数据，这可能与数据的精度与偶然性有关，也可能与模型参数估计偏于保守有关。

从拟合结果来看，男性和女性的功率—时间曲线趋势相同，功率随时间增加而呈 S 形下降。这与假设 2 是一致的。女性的 FTP 略小于男性，但其短时间内的极限功率高于男性。这体现了骑行者性别对该曲线的影响。其原因可能在于女性体重较轻，短时间内功率爆发较高，但长时间供氧能力较弱。两条曲线在经历一定时间之后的差异较小，这表明在相对较长的骑行时段内，性别所造成的差异可以忽略不计。

对于上述拟合结果，我们还从能量角度作了进一步的核对：由拟合所得的 $W'$ 与 FTP 可以反算出不同持续时间下的理论极限功率，其数值与由实测数据直接读取的极限功率相比，相对误差处于可接受的范围内；同时，拟合曲线在无氧供能向有氧供能过渡的临界时长附近并未出现突变，说明有氧池与无氧池的分段切换机制在实际数据上过渡平滑。这一核对在一定程度上支持了模型结构的合理性，也说明参数估计并未因分段处理而产生系统性的偏倚。

\subsubsection{不同类型骑行者的功率曲线}

除骑行者的性别之外，我们还应当考虑不同类型骑行者的特征，以确定其相应的功率—时间曲线。冲刺型选手（Sprinter）需要大量功率来对抗风阻，因此其体重与肌肉量较大，无氧池的能量储量大，但其功率体重比不太适合爬坡。为了对抗重力，爆发型选手（Puncheur）体重较轻，肌肉极限较低，冲刺能力较强，但不及冲刺型选手。全能型选手（Rouleur）各方面表现均衡，耐力强，通常在进攻时富有侵略性。为了直观地表现运动员的能力指标与模型参数之间的关系，各类骑行者指标的雷达图如下所示：

如上所述，选取 FTP、恢复系数 $k$、模型参数、体重、无氧池能量 $W'$ 作为指标，来表征各类骑行者的输出功率。根据各类骑行者的特点赋予不同的权重并进行归一化。在此基础上，我们为上述模型赋予不同的参数值，以得到相应的功率—时间曲线：

上图给出了各类骑行者的功率—时间曲线。由于各类骑行者能力不同，模型参数不同，曲线也因此不同。分析可得如下结论，这些结论与雷达图相一致：

\begin{itemize}
\item 全能型选手各方面均衡，因此其功率—时间曲线整体处于中间区段。其冲刺能力较强，同时也能在长时间骑行中维持一定的功率。
\item 冲刺型选手冲刺能力很强，其功率—时间曲线呈现出非常高的前段，即在短时间内极限功率很大，适合短时间的爆发式骑行。其无氧供能能力相对突出，无氧池能量储量大。
\item 爆发型选手耐力较差，其功率—时间曲线整体偏低，在短时间内具有一定的爆发力。其无氧供能能力略好。
\item 计时赛专家的综合能力最强，能够在长时间骑行过程中维持较高的功率。其有氧供能能力相对突出，FTP 较高，对应于上述模型中较大的有氧供能速率。
\end{itemize}

由上述结论可以看出，四类骑行者的差异可以被归结为模型参数在“短时项”与“长时项”上的不同配比：冲刺型选手的无氧池能量 $W'$ 与短时极限功率显著偏高，而 FTP 相对偏低，因此其曲线在前段陡峭、后段下沉；计时赛专家则反之，其 FTP 与有氧供能速率偏高，曲线整体抬高而在前段并不突出；爆发型选手介于二者之间但耐力最弱；全能型选手的各项指标居中，曲线因而最为平缓。这一结论说明，我们不必依赖繁琐的能力测试，只需由实测功率曲线拟合出少数几个参数，即可完成对骑行者的分类，并进一步以这些参数作为后续优化模型的输入。这正是疲劳生理模型与功率分配模型之间的接口所在。

\section{个体最优功率分配模型}

通过分析骑行者在骑行过程中的受力情况，并结合赛道条件，我们提出了个体最优功率分配模型，并使用动态规划方法寻找适合骑行的最优功率分配。

\subsection{数据处理}

我们需要将所建立的第一个模型应用于三条具体的赛道。因此，我们需要收集赛道的海拔数据与坡度（倾角）数据。相关的数据集与来源网站如下：

\begin{itemize}
\item 日本东京奥运会计时赛赛道的海拔与倾角数据，来源：\texttt{https://olympics.com/}；
\item 比利时佛兰德（Flanders）UCI 世界锦标赛计时赛赛道的海拔与倾角数据，来源：\texttt{https://www.flanders2021.com/}。
\end{itemize}

为了便于研究分析并将其应用于建模，我们以函数图像的形式给出所收集的数据。

由该图可知，三幅图中左上图表示东京奥运会计时赛的数据，右上图表示 UCI 计时赛的数据，下图为我们自行设计的赛道的数据。事实上，在 UCI 计时赛中，男子与女子选手的赛道不同。在奥运会计时赛中，男子与女子使用同一条赛道，但男子多骑行一圈。由于篇幅所限，具体情况将在结果分析中详细给出。

我们可以看到，UCI 计时赛赛道的海拔较低，较为平坦，但弯道略多。奥运会计时赛环线的赛道海拔较高，坡度较陡，弯道较少，且弯道集中分布在赛道的前半段。我们自行设计的赛道海拔介于两者之间，坡度集中在比赛的前半段与末端，共有五个弯道，且集中在比赛的后半段。我们设计的目的在于在一定程度上将坡度与弯道区分开来，以便于分别研究二者对骑行者功率的影响。

\subsection{模型的建立}

在建立模型的过程中，我们主要考虑骑行者骑行过程中所受到的阻力，分析骑行者的受力情况，并提出一个骑行阻力模型（Riding Resistance Model）。

\subsubsection{骑行阻力模型}

在个人比赛中，不考虑其他骑行者的干扰效应。因此，我们认为骑行者在骑行过程中所受到的阻力主要由两部分构成，一部分是空气动力阻力，另一部分是自行车与地面之间的滚动摩擦。根据他人的实验研究[8]，我们给出骑行者在平地运动过程中所受到的总阻力为：

\begin{equation}
F_{\mathrm{res}} = \frac{1}{2}\rho C_d A v_{ra}^2 + \mu m g,
\end{equation}

其中，$\rho$ 为流体密度，此处即空气密度；$A$ 为骑行者与自行车总投影的正对面积，即有效面积；$C_d$ 为阻力系数；$v_{ra}$ 为骑行者相对于流体的速度；$\mu$ 为滚动摩擦系数；$m$ 为骑行者与自行车的总质量；$g$ 为重力加速度。

该式的两项分别刻画了两类性质不同的阻力。第一项为空气动力阻力，它与速度的平方成正比，因而是高速骑行时的主要阻力来源，也是我们在优化中最需要关注的项——正是由于该项的非线性，功率分配才存在“何处加速、何处减速”的权衡；第二项为滚动摩擦阻力，它与速度无关而只与总重量及路面状况有关，因而在爬坡与低速段相对更为重要。

为了得到更为具体的表达式，我们对上述参数进行讨论，以便直接利用已有数据来简化求解。对于流体密度，我们有：

\begin{equation}
\rho = \rho_0 \cdot \frac{p}{p_0} \cdot \frac{T_0}{T},
\end{equation}

\begin{equation}
p = p_0\left(1 - \frac{0.0065\,h}{T_0}\right)^{5.256},
\end{equation}

其中，$\rho_0$ 为一个常数，一般取值为 1.293 kg/m$^3$；$p$ 为气压；$T$ 为以开尔文计的温度；$h$ 为当地海拔高度。

上述两式说明，空气密度并非一个可以随意取定的常数，而是随海拔升高而系统性地减小：海拔每升高一定高度，气压按幂律下降，密度随之减小，从而在相同速度下产生更小的空气阻力。这一点对于东京赛道这类海拔较高的赛道尤为重要，它解释了为何高海拔赛道上由空气阻力主导的那部分阻力会有所下降，同时也提示我们：重力分量的变化才是高海拔赛道的决定性因素。

对于投影有效面积，我们有：

\begin{equation}
A = k_A H^{0.725} M^{0.425}\left(1 - \varepsilon \sin\varphi\right),
\end{equation}

其中，$H$、$M$ 分别为骑行者的身高与体重；$\varphi$ 为骑行者头盔的水平偏转角，为简化问题，我们取其值为零。

对于阻力系数，我们有：

\begin{equation}
C_d = c_1 \mathrm{Re}^{c_2}, \qquad \mathrm{Re} = \frac{v_{ra} L}{\nu},
\end{equation}

其中，$\mathrm{Re}$ 为雷诺数，$L$ 为特征长度，$\nu$ 为空气的运动黏度，$c_1$、$c_2$ 为经验常数。

将阻力系数写成雷诺数的幂函数，是为了反映阻力系数并非固定不变这一事实：在自行车运动的实际速度范围内，绕流状态随速度变化，阻力系数随之略有变化。采用这一形式的优点在于，它使得阻力模型中唯一与速度强烈相关的部分仍然是 $v_{ra}^2$，而 $C_d$ 只提供温和的修正，从而在保证物理合理性的同时不显著增加求解的复杂度。

考虑滚动摩擦系数，根据经验数据[9]，山地自行车在沥青路面上的滚动摩擦系数为 0.003，而公路自行车的滚动摩擦系数通常更小。由于滚动摩擦约占总阻力的 10\%，并且为简化问题，我们不考虑赛道的材质，取 $\mu = 0.002$。

\subsubsection{赛道抽象模型}

对于个人自行车赛道，我们主要考虑骑行者在骑行过程中上下坡与转弯的情形分析。

对于上坡过程，我们将骑行者简化为一个质点并对其进行受力分析。

如上分析所示，我们可以简单地得到上坡过程的物理模型：

\begin{equation}
F = \frac{P}{v}, \qquad \frac{P}{v} = \frac{1}{2}\rho C_d A v_{ra}^2 + \mu m g \cos\theta + m g \sin\theta + m\frac{\mathrm{d}v}{\mathrm{d}t},
\end{equation}

其中，$F$ 为提供给骑行者的动力，它与骑行者所提供的功率有关；$\theta$ 为坡度角，当骑行者下坡时，它为钝角。

该式是骑行者在赛道上运动的牛顿第二定律表述：左端的 $P/v$ 表示把功率换算为牵引力（$P = Fv$），右端依次为空气阻力、滚动摩擦、重力沿坡面的分量与加速所需的惯性力。由于 $\theta$ 在平路上为零、上坡时为正、下坡时为钝角（其正弦为负），同一个式子即可统一描述上坡、平路与下坡三种情形，无需分段建立不同的方程，这也是我们采用带符号坡度角的动机之一。

需要指出的是，骑行者在下坡过程中的受力分析与上坡过程中的受力分析是相同的，上述物理模型对两种情形均适用。

对于弯道过程，为简化模型与计算，我们将该过程视为圆周运动。我们对其进行简单的受力分析。

根据上图，在弯道的切线方向上，骑行者受到总阻力以及自身所提供的动力。这与骑行者在直道上骑行时所做的相符。

在径向上，骑行者重力与路面支持力的合力提供向心力。然而，由于骑行者在转弯时相对路面的倾斜角有限，该合力是有限的。因此，在转弯过程中，骑行者的速度不能超过某一阈值。此外，我们根据转弯半径的大小将弯道分为两类：急弯与缓弯。

把上述径向条件写成不等式，即为
\[
\frac{m v^2}{R} \le m g \tan\beta ,
\]
其中 $R$ 为弯道半径，$\beta$ 为骑行者相对路面所允许的最大倾斜角。由此可得速度阈值
\[
v \le \sqrt{g R \tan\beta},
\]
即弯道半径越小、允许倾斜角越小，可通行的速度上限越低。这一不等式把“弯道”从几何信息转化为了对速度的硬约束，是后续动态规划中弯道处必须满足的可行性条件。

由此，我们得到对弯道过程的处理方式：

\begin{itemize}
\item 对于急弯过程，骑行者的速度不宜过高，骑行者应降低输出功率以将速度降至该阈值以下，否则无法完成转弯过程。也就是说，我们将弯道视为赛道上的速度检查点，只有低于速度阈值才能通过。
\item 对于缓弯过程，我们将其视为直道，不作特殊处理。
\end{itemize}

我们将先前得到的赛道高程图表示为海拔关于骑行者在赛道上位置 $x$ 的函数 $h(x)$。我们认为当海拔上升时骑行者处于上坡赛道，当海拔下降时处于下坡赛道。因此，通过对该函数求导，我们可以得到坡度角关于骑行者在赛道上位置 $x$ 的函数 $\theta(x)$。

同样地，由于我们简化了弯道模型，我们只考虑骑行者当前位置是直道还是弯道，而不关心弯道的半径。因此我们将赛道的倾角表示为一个二元函数 $I(x) \in \{0,1\}$，即：

\begin{equation}
I(x) = \sum_{i=1}^{n} \mathbf{1}\left[x \in \left[c_i - r_i,\; c_i + r_i\right]\right],
\end{equation}

其中，$c_i$ 为赛道中各个弯道的位置，$r_i$ 为该弯道的判定半宽，$n$ 为弯道总数；当 $I(x)=1$ 时表示该处为弯道，$I(x)=0$ 时表示该处为直道。

通过上述分析，我们将建模的基本思路总结如下。
\subsection{模型的求解}

我们需要利用已有的模型，在保证有氧池与无氧池的能量均不小于零的前提下，获得骑行者在一次比赛中的最优功率分配，以使骑行者能够在最短时间内到达终点。因此，采用动态规划算法求解该模型是一个良好的选择。

动态规划算法是一种将问题分解、并定义问题的状态与状态之间关系的算法，从而使问题能够以递推的方式求解。将该方法应用于上述模型时，主要思路如下：

\begin{itemize}
\item \textbf{划分子问题}。将赛道划分为一定数量的小段，并求出在每一小段具有相同的有氧剩余能量、无氧剩余能量以及末端速度时，完成这一部分赛道所需的最短时间。随着比赛进程逐步增加，直至覆盖整条赛道。那么完成整条赛道的最短时间，即为在任意有氧剩余能量、无氧剩余能量与末端速度条件下所计算出的最短时间。
\item \textbf{确定状态}。当有氧剩余能量为 $S_A$、无氧剩余能量为 $S_{An}$、末端速度为 $v$ 时，定义 $f_k(S_A, S_{An}, v)$ 为完成前方 $x_k$ 米赛道所消耗的最短时间。
\item \textbf{给定初值}。初始速度为零，有氧能量剩余量为 1，无氧能量剩余量为 1。
\item \textbf{状态转移}。对于每一小段，其状态函数的值可能来自赛道上一小段的所有状态值。因此，我们对赛道上一小段的所有状态值尝试所有功率。若其转移到自身的状态，则记录下来。遍历完成后取最小值。
\end{itemize}

相应地，状态转移方程可以写为：

\begin{equation}
f_k\left(S_A, S_{An}, v\right) = \min_{P \in \mathcal{P}} \left\{ \Delta t_k\left(P\right) + f_{k+1}\left(S_A', S_{An}', v'\right) \right\},
\end{equation}

其中，$\Delta t_k(P)$ 为在第 $k$ 小段上以功率 $P$ 骑行所耗费的时间，$\mathcal{P}$ 为允许的功率取值集合，$S_A'$、$S_{An}'$、$v'$ 分别为按前述消耗与恢复规律更新后的有氧剩余能量、无氧剩余能量与速度：

\begin{equation}
S_A' = S_A - r_A \Delta t_k, \qquad S_{An}' = S_{An} - r_{An}\Delta t_k \;\; \mbox{或} \;\; S_{An}' = S_{An} + k\frac{\mathrm{FTP} - P}{W'}\Delta t_k ,
\end{equation}

并要求 $S_A' \ge 0$、$S_{An}' \ge 0$。

需要说明的是，除时间之外，所有数值均被离散化以保证增量式遍历。当数值之间的差值足够小时，所得结果即可认为是准确的。

在具体实现中，我们将赛道按固定弧长划分为若干小段，并对有氧剩余比例、无氧剩余比例与末端速度三个连续量设置等间距的网格。这样做的好处是：状态的转移只可能发生在相邻网格点之间，因而遍历可以在有限步内完成；同时，由于我们以“到达某位置时的剩余能量与速度”为状态，而不以“已经用了多少时间”为状态，问题满足最优子结构，动态规划能够得到全局最优解而非局部最优解。计算中，若某个状态组合所对应的剩余能量不足以支持该小段的任何可行功率，则该状态被标记为不可行并从后续转移中剔除；这一可行性判定等价于在优化中加入能量的硬约束，从而保证所求得的解在生理上是可以真正执行的。当网格加密到一定程度后，继续加密所引起的完赛时间变化已小于秒级，说明离散化误差已被控制在可接受的范围内。

\subsection{结果与分析}

考虑到骑行者的性别与类型，在将模型应用于赛道时，我们分别讨论了分为男性和女性的计时赛专家与冲刺型选手。

在三条赛道中，两类男子骑行者的功率分配策略如下所示。图中纵向依次为冲刺型选手与计时赛专家，横向依次为东京赛道、比利时赛道与我们自行设计的赛道。

阅读这些图像时应注意两点。其一，横坐标为比赛进程（即骑行者在赛道上的位置），纵坐标为按体重归一化后的输出功率，因而曲线的上下移动表示输出强度的变化，而曲线的宽窄并不直接对应路段的长度；其二，同一条曲线上相邻的峰与谷之间并非彼此独立，一次冲刺所透支的无氧储备必须在随后的低功率区间内被偿还，因此曲线的“峰高”与“谷宽”总是成对出现，二者的比例由 (4) 式所描述的恢复速率决定。正是这一约束使得最优策略无法通过简单地“处处全力”来获得，而必须在赛道上反复进行取舍。

基于上图，我们得到了不同类型骑行者对应的最优功率分配。按照该策略，在比赛进程的相应位置施加相应的功率，骑行者将在自身条件下以最短的时间完成比赛。当曲线急剧上升时，骑行者应当发起冲刺；当曲线急剧下降时，骑行者应当降低功率以恢复无氧池的能量；在大部分时间内，曲线围绕 FTP 上下波动，骑行者应保持相对稳定的功率骑行。我们还绘制了无氧池能量与有氧池能量的变化以佐证这一点，详情可参见附录。此外，我们可以得出以下结论：
\begin{itemize}
\item \textbf{各类骑行者在同一赛道上的功率分配策略基本相同。}由前述内容可知，不同类型骑行者的功率曲线具有相似的趋势。当它们被应用于相同的赛道条件时，所得到的功率分配策略也具有相同的趋势。
\item \textbf{骑行者在高海拔比赛中需要更频繁地冲刺。}虽然海拔较高处的平均阻力小于海拔较低处，但坡度变化更大，因此骑行者需要提高功率以爬坡，随后降低功率以恢复。所以，在起伏更为密集的丘陵路段上，骑行者的功率波动更大。这也同时说明，重力对骑行者的影响大于阻力的影响。
\item \textbf{弯道会降低骑行者的功率输出。}从自行设计的赛道可以看出，骑行者在弯道处存在一定程度的功率降低。这也与我们此前的分析相一致。
\item \textbf{冲刺型选手在赛程前半段可以发起比计时赛专家更强的进攻。}横向比较上述图像，在比赛的前半段，可以看出冲刺型选手发起冲刺所对应的功率峰值更高。在比赛的后半段，由于冲刺型选手的耐力略差，冲刺次数与功率峰值略有减少。
\item \textbf{冲刺型选手受弯道的影响小于计时赛专家。}在弯道更为密集的东京赛道上，冲刺型选手的功率削减次数更少。这可能与该类骑行者较低的平均功率有关。
\item \textbf{冲刺型选手受海拔的影响大于计时赛专家。}在海拔更高、坡度变化更陡的奥运赛道，冲刺型选手的功率曲线波动中心更低，平均功率也更低。
\end{itemize}

上述策略的波形结构并非偶然，它可以由模型的约束条件逐项解释：在坡道之前，坡度的上升会立即抬高重力分量，若不在坡底提前积蓄速度，则骑行者必须在坡上以更高的功率维持速度，而在坡上提高功率的效率远低于在坡底提高功率，因此最优解总是把冲刺安排在坡底；在下坡处，重力分量为负，骑行者即使不输出功率也会加速，此时把功率降到 FTP 以下不仅不会损失时间，反而可以恢复无氧池的能量，为下一次冲刺做准备；而在平路上，空气阻力与速度的平方成正比，超过某一速度之后边际收益迅速衰减，因而最优解将功率稳定在 FTP 附近，使有氧供能持续而不过早透支。换言之，最优功率曲线是“能量约束下的边际收益均衡”这一原则在赛道上的逐点体现。

由于篇幅所限，女子比赛的最优功率策略图与女子 UCI 赛道图附于附录之中。由于奥运会赛道与 UCI 赛道中男子与女子赛道不同，为控制变量，我们仅对自行设计的赛道进行分析，并可得出以下结论：

\begin{itemize}
\item 女子的平均功率低于男子。可以看出女子的功率曲线整体偏低。
\item 女子在比赛前半段的冲刺强度可以达到更高，甚至高于男子。女子功率曲线在比赛前半段的峰值更高。
\item 女子受坡度与弯道的影响大于男子。女子功率图在连续的坡道与弯道处波动更为频繁。
\end{itemize}

上述性别差异同样可以从模型参数中得到解释：女子的体重较轻，因而在相同坡度下所需克服的重力分量较小，但按体重归一化后的无氧储量 $W'$ 相对有限，故其可用于冲刺的“预算”较小，一旦遇到连续的坡道与弯道，就需要更频繁地在冲刺与恢复之间切换，功率曲线因而显得更为“毛糙”。这也意味着对于女子比赛而言，赛前的功率分配预案应当预留更大的裕度，并在执行时更强调对节奏的控制。

此外，我们给出了骑行者所需最短时间的统计图，为分析提供了直观的数据。

通过该图，我们给出了不同骑行者在最优功率分配下的具体用时。如上所述，对于男子与女子之间的比较，我们仅以自行设计的赛道为参考。可以看出，男子选手的用时更短，领先女子选手 172 s，且平均功率更高。此外，计时赛专家的总体用时短于冲刺型选手，平均领先后者约 300 s。他们的平均功率更高，综合能力更强，且受海拔的影响更小。特别是在高海拔赛道上，前者领先后者 371 s，这是三条赛道中两者差距最大的一条。这与我们先前的模型高度吻合，同时也验证了模型的有效性与合理性。

对这些数字还可以作进一步的解读。计时赛专家相对于冲刺型选手的优势随赛道海拔的升高而扩大（在自行设计赛道上约为 300 s，在奥运赛道上达到 371 s），这一趋势说明：赛道越陡，有氧能力在成绩中的权重就越大，而无氧爆发能力的相对价值则被削弱。这是因为在连续爬坡的路段上，骑行者的功率长期处于 FTP 附近而无法依赖短时爆发，谁能以更高的有氧功率长时间输出，谁就能在整条赛道上积累优势；相反，在平坦赛道上空气阻力占主导，冲刺能力可以在有限的几个节点上被有效利用，两类选手的差距因而缩小。男子与女子之间 172 s 的差距同样可以由平均功率的差异来解释，且与平均功率之比基本相当，说明模型给出的差异主要来源于能力参数而非数值处理的偶然性。

\section{灵敏度分析}

\subsection{天气与环境影响}

为考虑天气的潜在影响，我们对模型进行了修正。我们考虑了风向与风力（即风速）对模型参数的影响，并作出新的假设：

风向与风速只是与空间有关的函数，在空间中形成一个风的矢量场，而与时间无关。并且我们只把风视为水平方向的气流。这是由于赛场上风速与风向的变化缓慢，而且通常只检测一小时内的平均风速。我们所采用的风速与风向数据如下：东京赛道男子比赛的风向为 $270^\circ$、风速为 0.5 m/s；东京赛道女子比赛的风向为 $290^\circ$、风速为 1.0 m/s；比利时赛道男子比赛的风向为 $050^\circ$、风速为 2.0 m/s；比利时赛道女子比赛的风向为 $010^\circ$、风速为 6.0 m/s[10]。

风对模型结果的影响只通过空气阻力产生。为简化模型，我们只考虑与骑行者行进方向相同或相反的风对骑行产生影响，垂直方向的风不影响骑行。由此，我们将先前模型中的总阻力修正为：

\begin{equation}
F_{\mathrm{res}} = \frac{1}{2}\rho C_d A\, v_{ra}^2 \, \mathrm{sgn}\left(v_{ra}\right) + \mu m g, \qquad v_{ra} = v + v_w \cos\varphi ,
\end{equation}

其中，$\varphi$ 为行进方向与风向之间的夹角（小于或等于平角），它本质上是 $x$ 的函数，记为 $\varphi(x)$；$\mathrm{sgn}(\cdot)$ 为符号函数；$v$、$v_w$ 分别为骑行者的水平速度与风速。

该修正的实质是把空气阻力的表达式中的相对速度由 $v$ 替换为 $v_{ra} = v + v_w\cos\varphi$：当风为顺风时 $\varphi = 0$，相对速度增大，阻力随之增大，但骑行者实际获得的地速更高；当风为逆风时 $\varphi = \pi$，相对速度被进一步放大，阻力迅速增大；当侧风时 $\varphi = \pm\pi/2$，按本模型的简化假定，其影响为零。保留符号函数是为了处理极端情形：当逆风风速超过骑行者地速时，$v_{ra}$ 变为负值，此时阻力方向反转，符号函数保证了力的方向仍然正确，从而避免模型在这一退化情形下给出无意义的结果。

综上所述，新模型是对原模型的改进。我们可以利用上述数据与类似的方法求解该模型，并通过改变风向与风速的取值来分析灵敏度。

为了探究天气条件对结果的影响，我们以女子在 UCI 赛道上的表现为例，设置了 8 个风向与 3 个风力等级，共 24 种风环境。在这些环境中，我们模拟了运动员在大风中的表现。所得到的结果如下所示。

风向采用正北、正南、正西、正东、东北、东南、西北、西南 8 个梯度设置；风速采用 1 m/s、2 m/s 与 3 m/s 三个梯度设置。

需要说明的是，之所以选取 1$\sim$3 m/s 作为风速的取值范围，是因为该范围覆盖了比赛日实测风速的常见区间，且在这一量级下，风对完赛时间的影响已足以显现而不会使骑行者无法完成比赛；之所以选取 8 个风向而非连续变化，是为了使结果可以被直观地比较与解释。24 种组合的模拟均沿用此前的最优功率分配求解流程，唯一的差别在于每一小段上的阻力按修正后的 (14) 式计算，因而所得到的完赛时间差异可以完全归因于风的作用。

根据上图，当风向一定时，风速每变化一个等级，完赛时间将变化 20$\sim$330 s，变化比例小于 10\%。当风力一定时，风向每向一个方向变化，完赛时间将变化 30$\sim$460 s，最大变化比例为 16\%。因此，综合考虑两个因素的影响，我们认为我们的模型对天气环境的变化较为敏感。

进一步地，从 24 种环境的分布可以看出，顺风与逆风之间的差异是最主要的来源：逆风情形下完赛时间显著延长，顺风情形下则相应缩短，而侧风情形基本落在两者之间。这一结果与 (14) 式的结构完全一致，也从侧面说明模型对风向的处理是自洽的。此外，由于风对结果的影响主要体现在空气阻力项上，而该项在高速段占主导，因此风的影响在平坦、速度较高的赛段上更为显著，而在爬坡、速度较低的赛段上相对较小。这一结论对实际比赛具有直接的指导意义：在逆风条件下，骑行者应避免在平路高速段过度消耗无氧储备，而应把冲刺留给坡度较大的路段。

\subsection{偏差影响}

在比赛中，骑行者无法精确地执行所确定的最优功率分配策略。例如，骑行者在本应提高功率的赛段上并未提高到最优功率。因此，骑行者的功率—位置曲线很可能偏离最优曲线。为模拟这一偏差，我们对最优功率分配曲线进行随机扰动，使骑行者可能在任意位置输出次优功率。我们以骑行者实际输出功率曲线与最优功率曲线在所有位置上偏差绝对值的平均值，作为骑行者偏离最优曲线的程度。

在具体操作上，我们以最优功率曲线 $P^{\ast}(x)$ 为基准，对赛道上的每一点施加一个随机扰动 $\delta(x)$，得到实际执行的功率曲线 $P(x) = P^{\ast}(x) + \delta(x)$，并通过调整扰动的幅度使平均偏差 $\overline{|\delta|}$ 达到预设的若干个水平。需要强调的是，扰动后的曲线同样要满足能量约束：若扰动导致某一位置后的剩余能量为负，则该组扰动被视为不可行而重新抽样。这样得到的每一组“非最优但可行”的功率分配，才是现实中一位骑行者可能真正执行出来的功率分配。

我们使用随机数生成非最优的功率分配，并将其代入模型求解相应的完赛时间，结果如下：当平均偏差为 0、0.06、0.07、0.08、0.09、0.1、0.15、0.2（单位 W/kg）时，完赛时间分别为 4001.5 s、4010.8 s、4144.1 s、4159.9 s、4174.4 s、4196.7 s、4200.0 s、4325.7 s；相对于最优解的用时增量分别为 0 s、$+9.3$ s、$+142.6$ s、$+158.4$ s、$+172.9$ s、$+195.2$ s、$+198.5$ s、$+324.2$ s；相应的变化比例分别为 0、0.232\%、3.563\%、3.959\%、4.142\%、4.878\%、4.961\%、8.101\%。

由表中数据可以发现，随着偏差值的增大，骑行者的完赛时间逐渐增加，这在一定程度上验证了我们的模型确实是最优解。当平均偏差低于 0.06 时，完赛时间的变化很小，小于 0.5\%。当偏差增大到 0.06 与 0.07 之间的某一数值之后，完赛时间的变化增大，但仍小于 5\%。而当平均偏差处于 0.15 与 0.2 之间时，完赛时间的变化大幅增加，超过 5\%。若平均偏差继续增大，骑行者将由于无法满足有限的能量约束而无法正常完成比赛。

经过上述分析，当平均偏差变化约 5\% 时，完赛时间变化约 10\%。因此，我们认为模型结果对骑行者偏离目标功率这一情形具有一定的不敏感性。

这一不敏感性具有明确的机理：最优解之所以在功率曲线上表现为围绕 FTP 的波动，是因为在平路段，完赛时间对功率的响应是凸的但在最优功率附近较为平坦，功率的微小偏离只会带来二阶小量的时间损失；而在坡道与弯道处，虽然功率的时间灵敏度较高，但这些路段在整条赛道中所占的比例有限，因而其对总时间的影响被稀释。换言之，模型的“宽容”并不意味着功率分配无关紧要，而是意味着骑行者无须以机械的方式逐点复现最优曲线，只要把握住关键节点的节奏，即可获得接近最优的成绩。

此外，我们将整条赛道划分为若干段，每一段的端点即为一个检查点。通过多次模拟，我们得到了到达各个检查点的时间区间。若骑行者在这些区间内到达相应的检查点，则最终的完赛时间可以保证不晚于最优完赛时间之后 1 分钟。下表给出了将 UCI 男子个人计时赛赛道划分为 11 段之后，到达各个检查点可接受的时间范围：在 4、8、12、16、20、24、28、32、36、40 km 处，可接受的最晚到达时间分别为 287、717、1188、1469、1856、2193、2518、3035、3316、3710 s，最早到达时间分别为 283、703、1148、1429、1815、2152、2477、2974、3274、3665 s。

这一结果的实际价值在于把“最优功率曲线”转化为了“可执行的时间表”。在比赛中，选手无法时刻盯住功率计上的数字，也难以判断自己当前是否偏离了最优曲线，但却可以通过路标与码表方便地判断自己到达某一距离的时间是否落在允许区间内。区间宽度在不同路段并不相同：在坡度变化剧烈、弯道密集的路段，区间相对较窄，说明这些路段对节奏更为敏感，选手须格外留意；而在平坦的长直道上，区间相对较宽，选手可以有一定的调整余地。若某一检查点的到达时间明显早于最早时间，说明前半程过度消耗了无氧储备，应当立即降低功率；若明显晚于最晚时间，则说明前面的节奏过于保守，可在剩余赛程中适度加大输出。

\section{模型评价与进一步讨论}

\subsection{优点}

我们的模型具有以下优点：

\textbf{我们对三罐模型进行了更为深入的分析，并提出了有氧与无氧能量疲劳模型。}借鉴他人的研究，我们将骑行者的功率输出划分为有氧部分与无氧部分。最终，我们发现我们的模型对数据拟合得相当好，模型的有效性得到了验证。

这一划分的意义不仅在于拟合优度，更在于它使功率曲线从一条经验曲线变成了可以被逐段解释的物理对象：曲线的短时端由无氧储量与消耗速率决定，长时端由有氧功率决定，而二者的衔接则决定了曲线的“肩部”位置。正因为每个参数都有明确的生理含义，我们才能据其完成骑行者的分类，并把参数直接传递给后续的优化模型。

\textbf{综合考虑了多种因素，提出了骑行者赛道模型。}在这一过程中，我们考虑了骑行者的能量消耗极限与受力情况、赛道的坡度与倾角、风向与风速等多种因素，因而更为全面。

与只考虑单一阻力或单一坡度的简化模型相比，我们的模型在同一框架内同时容纳了生理约束与几何约束，因而能够回答诸如“高海拔赛道上应当更频繁地冲刺还是更均匀地输出”这类需要多因素权衡的问题，而这些结论正是单纯依赖经验配速所无法得到的。

\textbf{模型具有很强的可推广性。}在建模时，我们从个人计时赛入手。随后，我们根据团队比赛的特点对模型进行了修正，并提出了轮换领骑策略，从而将模型推广到适用于团队计时赛的情形。

推广过程中被修改的仅仅是阻力模型中的一个系数，而能量模型与优化框架均原封不动地保留下来。这说明我们所建立的模型在结构上具有模块化的优点：当比赛规则或环境条件发生变化时，只需替换相应的模块，而不必重建整个模型。

\subsection{缺点}

我们的模型存在以下局限及相应的改进方向：

\textbf{在无氧能量疲劳模型中，我们将恢复系数视为常数。}这样做简化了模型，但在一定程度上降低了模型的真实性。事实上，由于能量的损耗，将恢复系数描述为变量更为合适。

一个可行的改进是把 $k$ 写成剩余能量的函数 $k(S_{An})$：当无氧池接近耗尽时，恢复过程应当更为缓慢；当无氧池较为充盈时，恢复则可以更快。另一种改进是引入恢复的“滞后效应”，即恢复速率不仅取决于当前的功率缺口，还取决于此前透支的深度与持续时间。这类改进不会改变模型的基本结构，只会增加一个待辨识的函数。

\textbf{在构建个体最优功率分配模型时，对赛道抽象模型的考虑相对简单。}我们只关心比赛进程的某一位置是否存在弯道，而未考虑弯道半径的大小。虽然模型得以简化，但模型的真实性被削弱。事实上，这一因素可能影响骑行者的低功率骑行时间，进而影响结果。

若进一步把弯道半径纳入模型，则弯道不再只是“是否”的二元标记，而会给出一个与半径相关的速度上限，从而影响骑行者在弯道上的减速幅度与出弯加速的时机。对于半径差异较大的赛道（例如既有高速大半径弯道又有城市街区内的直角急弯），这种改进的必要性尤为明显。

\subsection{模型的推广}

本节我们主要讨论模型的推广。我们此前所建立模型的背景是个人计时赛。至于团队计时赛，亚里士多德曾说：“整体大于部分之和。”因此，我们需要进一步修正模型。

根据 UCI 给出的团队计时赛规则[11]，不同队伍之间应保持一定的距离，队员之间不得相互推行，也不得使用跟随的辅助车辆。因此，队伍之间的影响以及其他影响阻力、动力与能量消耗的因素均被忽略。重要的是，在团队计时赛中，前排选手可以通过前后排列为后排选手减小空气阻力。因此，团队计时赛与个人计时赛的主要区别在于，我们需要考虑不同骑行者之间阻力的差异。

从另一个角度看，这一差异把原本彼此独立的个体优化问题变成了一个耦合问题：每名队员的功率输出会影响队形的位置关系，而队形的位置关系又反过来改变各队员所受的阻力。正因如此，团队计时赛中的“最优”不再是单条功率曲线，而是一组功率曲线与一组换位时机的组合。

阻力模型首先需要修正如下：

\begin{equation}
F_{\mathrm{res},i} = \eta_i \cdot \frac{1}{2}\rho C_d A\, v_{ra}^2 + \mu m g ,
\end{equation}

其中，$\eta_i$ 为一个常数，与骑行者总人数以及该骑行者在队列中的位置有关。

当队伍总人数为 6 时，各个位置上骑行者所受阻力表达式中的 $\eta_i$ 取值如下：第一至第六位依次为 95\%、68\%、67\%、67\%、67\%、69\%。

此外，我们需要考虑的决策不仅限于何时加速与何时减速。为了充分利用队伍带来的减阻效果，团队作为一个整体需要保持相对静止的队形，并以一定的队伍速度前进。由于队伍后部的选手所面对的阻力更小，运动更为省力，而前排选手则更为费力。

在团队计时赛中，队伍往往采用“轮换领骑”（Draft）来实现队伍工作负荷的平衡。“轮换领骑”的过程是：队伍的第一位选手横向移出队列，使第二位选手成为队伍的领骑者，随后第一位选手减速，移动到队伍末尾。该过程的示意图（以四名选手为例）如下所示。

因此，何时、在何地“轮换领骑”是团队计时赛中的一项关键决策。于是，我们定义一个离散取值函数来表示替换情形：

\begin{equation}
D(x) = \begin{cases} 1, & x = x_j, \quad j = 1,2,\dots,J, \\[2pt] 0, & \mbox{其他}, \end{cases}
\end{equation}

其中，$x_j$ 为比赛进程中需要轮换领骑的位置，$J$ 为轮换次数。

该函数提供了一个新的决策变量，我们可以通过更为复杂的动态规划来寻找最优解。

此外，由于比赛成绩以第四名队员的成绩计，我们可以结合不同队员的不同能力，在行进过程中的某些位置放弃较弱的选手，以获得更好的成绩。

这一规则带来了一个与个人计时赛本质不同的自由度：既然只有第四名队员的成绩被计入，那么队伍就不必强求所有队员同时到达，而可以在能力较弱的队员完成其应尽的领骑任务之后让其掉队。相应地，队形的人数会在比赛中途发生变化，而人数的变化又会改变剩余队员所受的阻力系数，因此“何时掉队”与“何时轮换”必须被联合决定。

团队计时赛中还有更多的策略。为突出目标，我们仅选取一种有代表性的策略以简化决策。我们也可以考虑这样的情形：原本处于队伍前部的选手不移动到队尾，而是与第二位选手交换位置，或者回到队尾之外的其它位置。这同样可以根据不同骑行者能力水平的不同进行修改。

在团队计时赛中，基于我们此前的模型，我们可以考虑更多的约束条件与决策变量，并提出改进的模型以用更为复杂的动态规划求解。由于篇幅所限，此处不再描述模型的求解过程。

\section{结论}

本文建立了一个适用于各类骑行者的模型，用以分析骑行者应施加的功率与比赛位置之间的关系，并给出最优的功率分配策略。

首先，我们建立了有氧与无氧能量疲劳模型。在此基础上，我们定量地给出了不同骑行者类型与性别的功率曲线，并使用模型参数来定义不同类型与性别的自行车运动员。

其次，我们考虑了骑行者的受力、赛道坡度与转弯情形，并结合运动员的功率曲线建立了最优个体功率分配模型。我们使用动态规划得到了不同自行车运动员的最优功率分配图。此外，我们还给出了不同运动员对赛道环境（如坡度与弯道）适应性的特征，并计算了他们的最短用时。

第三，我们考虑了风向与风速对模型作出修正，并针对天气条件进行了灵敏度分析。我们发现，风速变化 1 m/s 会使结果变化不超过 10\%，风向变化 45 度会使结果变化不超过 16\%。因此，我们认为我们的模型对风向与风速较为敏感。我们还考察了模型对功率分配偏差的灵敏度，结果发现我们的模型对偏差并不敏感。随后我们给出了可接受的偏差范围。

最后，我们考虑了风阻对团队的影响，并将个人计时赛模型推广为团队计时赛模型。我们给出了关于轮换领骑与掉队策略的考虑。并在文末为骑行者和车队经理（Directeur）提供了一份比赛指南。

综合全文，可以得到这样一条贯穿始终的主线：骑行者的成绩取决于他如何在“能输出多大功率”与“能把功率维持多久”之间作出取舍，而这种取舍在赛道上的每一个位置都以不同的形式出现——在坡道上体现为对抗重力，在弯道上体现为对抗速度上限，在平路上体现为对抗空气阻力，在有风时体现为对抗相对速度的变化。我们的模型所做的，就是把这条主线转化为一个可以在计算机上求解的优化问题，并把求解结果翻译为骑行者可以直接执行的功率—位置曲线与时间检查点。模型在灵敏度分析中所表现出的“对天气敏感、对偏差不敏感”的特性，也恰好对应了实际比赛中教练与选手所关心的两个问题：赛前的天气评估值得投入精力，而赛中的节奏执行则可以留有余地。

\section*{男子精英个人计时赛参赛指导}

\noindent\emph{——2021 年比利时佛兰德 UCI 世界锦标赛}

\subsection*{赛道与天气分析}

比赛在比利时佛兰德省举行。骑行者需要从克诺克—海斯特（Knokke Heist）骑行至比利时北部的布鲁日（Bruges）。

赛道的高度波动情况如下。赛道总长 43.3 km，平均海拔 15 m，极差为 15 m，标准差仅为 3 m。赛道平坦，更适合平路型计时赛专家（Flat Time Trial Specialist）。

赛道共有 7 个急弯，分别位于 1.5 km、2.4 km、2.7 km、11.2 km、14.3 km、17.6 km 与 27.6 km 处。注意：骑行者经过这些弯道时需要提前减速，建议将速度控制在 10 m/s 以下，以免冲出赛道。

值得注意的是，前三个急弯彼此相距很近，集中出现在赛程的开头 3 km 之内。这意味着选手在出发后不久就要连续三次削减功率，若在此处过早发力，将白白消耗无氧储备而在随后的长直道上无从恢复。因此，建议把开局的节奏控制得更保守一些，等通过 2.7 km 的急弯之后再回到比赛节奏。后半程的急弯则分别出现在 11.2 km、14.3 km、17.6 km 与 27.6 km 处，间隔较为均匀，可以作为天然的节奏节点使用。

根据美国国家海洋和大气管理局（NOAA）所提供的泽布吕赫（Zeebrugge）气象观测站数据，在比赛当日，即 2021 年 9 月 19 日 16:00，风向为北偏东 $50^\circ$，风速为 2 m/s，属于顺风。风对成绩的影响约为 1\%，约 $\pm 10$ s。

\subsection*{自行车运动员特征分析}

根据所给定的骑行者数据，可以得到该骑行者的功率曲线。其 FTP 为 4.79 W/kg，因此其有氧呼吸能够提供较高的功率效率。其无氧总能量为 56.28 J/kg，因此其无氧呼吸的短时功率效率较差。这意味着他的耐力较好，而短期爆发力较差。

他是典型的平路型计时赛专家。他更适合在丘陵较少的平坦地形上骑行，这与该赛道的特点相符。

这一功率画像对比赛策略的直接含义是：他应当以较高的稳定功率贯穿全程，而不是依赖多次短促的爆发。换言之，他的最优策略更接近于一条在 FTP 上下小幅波动的平滑曲线，冲刺的次数应当少而有节制，且优先安排在急弯出弯后的短暂加速段，以尽快回到巡航速度。若赛道上出现连续的短坡，他应更早地降低功率以保存无氧储备，因为其无氧总能量偏小，一旦透支便难以在剩余赛程中恢复。

\subsection*{行进方向概览}

通过综合分析赛道模型与选手模型，对于这位希望取得最佳成绩的骑行者，我们可以得到推荐的功率—距离分布图（P-D 图）。

可以看出，存在若干需要施加更大力量以加速的关键节点。此外，也存在需要降低功率以恢复体力的节点，包括在下坡位置以及急弯之前降低功率。建议骑行者在 $x = 1$ km、9.6 km、17.8 km 等位置提高输出功率以完成加速；而在 $x = 1.6$ km、2.7 km 等位置将功率输出降低至约 1 W/kg。其余时间应尽量维持在 FTP 附近，即有氧呼吸的最大值。

按照该曲线完成整条赛道可以使选手的表现达到最优，预计完赛时间为 3499 s（58 分 19 秒）。

然而，高度细致地遵循上述功率分配是不可能的。功率输出上的小幅偏差仍可带来接近最优的结果。因此，我们为骑行者到达某些位置提供了可接受的时间区间，其中一部分如下所示。只要选手到达这些位置的时间处于最晚时间与最早时间之间，最终完赛时间就可以不晚于 57 分 19 秒，即比最优时间晚 1 分钟。

在 4、8、12、16、20、24、28、32、36、40 km 处，可接受的最晚到达时间分别为 287、717、1188、1469、1856、2193、2518、3035、3316、3710 s，最早到达时间分别为 283、703、1148、1429、1815、2152、2477、2974、3274、3665 s。

\subsection*{模型总结}

我们的模型从研究骑行者的有氧呼吸与无氧呼吸入手，确定了能量消耗与功率输出之间的关系。该模型刻画了骑行者的能量特征，并在外部表现为功率曲线。对于任意给定的赛道，我们模拟了骑行者状态与赛道上各种阻力之间的关系。在结合上述模型时，我们使用动态规划得到了最优输出功率与赛道位置之间的关系，这对骑行者具有很大的参考价值。只要骑行者按照给定的功率输出进行比赛，就能取得最佳成绩。为了降低选手执行的难度，我们提供了相对宽松的参考时间，同样可以获得接近最佳的成绩。

对车队经理而言，这份指南还提示了三项可以提前准备的工作：其一，赛前应当根据当日实测的风向与风速重新运行一次模型，因为结果对风的敏感程度远高于对执行偏差的敏感程度；其二，应当把检查点的时间区间打印在车把上的提示卡上，使选手无需计算即可判断自己的节奏是否偏离；其三，应当根据赛道的坡度分布与急弯位置预先划定“冲刺段”与“恢复段”，并在赛前训练中反复演练，使选手在比赛中能够凭感觉复现这些节点，而不必依赖实时的功率读数。

\section*{参考文献}

\begin{itemize}
\item[{[1]}] Poole, David C et al. ``Critical Power: An Important Fatigue Threshold in Exercise Physiology.'' \emph{Medicine and science in sports and exercise} vol. 48,11 (2016): 2320--2334. doi:10.1249/MSS.0000000000000939
\item[{[2]}] \texttt{https://commons.wikimedia.org/w/index.php?search=road+cycling}
\item[{[3]}] Jones, A.M., Vanhatalo, A. The `Critical Power' Concept: Applications to Sports Performance with a Focus on Intermittent High-Intensity Exercise. \emph{Sports Med} 47, 65--78 (2017).
\item[{[4]}] Morton RH. A 3-parameter critical power model. \emph{Ergonomics}. 1996 Apr;39(4):611--619. DOI: 10.1080/00140139608964484. PMID: 8854981.
\item[{[5]}] Coggan, Andrew R. ``Training and racing using a power meter: an introduction.'' (2003).
\item[{[6]}] Monod, H., \& Scherrer, J. (1965). The work capacity of a synergic muscular group. \emph{Ergonomics}, 8(3), 329--338.
\item[{[7]}] Morton, R. H. (1986). A three component model of human bioenergetics. \emph{Journal of mathematical biology}, 24(4), 451--466.
\item[{[8]}] Crouch, T.N., Burton, D., LaBry, Z.A. et al. Riding against the wind: a review of competition cycling aerodynamics. \emph{Sports Eng} 20, 81--110 (2017). \texttt{https://doi.org/10.1007/s12283-017-0234-1}
\item[{[9]}] Pierre Debraux, Frederic Grappe, Aneliya V. Manolova \& William Bertucci (2011) Aerodynamic drag in cycling: methods of assessment, \emph{Sports Biomechanics}, 10:3, 197--218, DOI: 10.1080/14763141.2011.592209
\item[{[10]}] Steyn, Wynand J. vdM, and Janike Warnich. ``Comparison of tyre rolling resistance for different mountain bike tyre diameters and surface conditions.'' \emph{South African Journal for Research in Sport, Physical Education and Recreation} 36.2 (2014): 179--193.
\item[{[11]}] \texttt{https://www.ncei.noaa.gov/products/land-based-station/integrated-surface-data-base}
\end{itemize}

\section*{附录}

附录 1 为女子选手的分析图集，其中包括：UCI 女子赛道图；女子选手在奥运会赛道上的功率输出；女子选手在 UCI 赛道上的功率输出；女子选手在自行设计赛道上的功率输出；以及女子选手有氧与无氧剩余能量的变化曲线。这些图件用于佐证正文第 5.4 节中关于女子选手功率分配策略与能量变化的结论。

\end{document}
